% !TeX program = pdflatex
\documentclass{beamer}
\usepackage[utf8]{inputenc}
\usepackage[polish]{babel}
\usepackage[T1]{fontenc}
\usepackage{tikz}
\usetikzlibrary{positioning, arrows.meta, decorations.pathreplacing, calc}
\usepackage{graphicx}
\usepackage{hyperref}
\usepackage{biblatex}
\usepackage{amsmath, amssymb}
\usepackage{caption}
\addbibresource{references.bib}

% Global font settings
\usepackage{lmodern}  % Latin Modern font
\usefonttheme{professionalfonts}  % Use professional fonts
\setbeamerfont{title}{size=\large,series=\bfseries}
\setbeamerfont{frametitle}{size=\normalsize,series=\bfseries}
\setbeamerfont{normal text}{size=\small}

% Theme and colors
\usetheme{Madrid}
\usecolortheme{dolphin}
\setbeamertemplate{navigation symbols}{}
\setbeamertemplate{footline}[frame number]

% Diagram styles
\tikzset{
  block/.style={
    draw,
    rectangle,
    rounded corners,
    minimum width=2cm,
    minimum height=1cm,
    align=center,
    fill=blue!10
  },
  annotation/.style={
    draw=red,
    thick,
    -Stealth
  }
}

% Formula environment
\newenvironment{formula}{%
  \begin{tcolorbox}[colback=yellow!10!white, colframe=black!60, boxrule=0.5pt, arc=3mm, boxsep=3pt]%
}{%
  \end{tcolorbox}%
}

% Add tcolorbox package
\usepackage{tcolorbox}

\title{Podsumowanie wyników modeli generatywnych}
\subtitle{Analiza porównawcza różnych architektur}
\author{Jakub Barylak, Michał Jagoda, Piotr Marcol}
\institute{Politechnika Śląska}
\date{Maj 2025}

\begin{document}

\begin{frame}
  \titlepage
\end{frame}

\begin{frame}{Plan prezentacji}
  \tableofcontents
\end{frame}

\section{Przegląd zaimplementowanych modeli}

\begin{frame}{Zaimplementowane modele}
  \begin{columns}[T]
    \begin{column}{0.5\textwidth}
      \textbf{Podstawowe modele:}
      \begin{itemize}
        \item Klasyczny Autoencoder
        \item Wariacyjny Autoencoder (VAE)
        \item Conditional VAE
      \end{itemize}
    \end{column}
    \begin{column}{0.5\textwidth}
      \textbf{Zaawansowane modele:}
      \begin{itemize}
        \item Generative Adversarial Network (GAN)
        \item Vector Quantized VAE (VQ-VAE)
        \item Diffusion Model
      \end{itemize}
    \end{column}
  \end{columns}
\end{frame}

\section{Wyniki i metryki}

\begin{frame}{Metryki oceny}
  \begin{block}{Wykorzystane metryki}
    \begin{itemize}
      \item \textbf{Basic AE:} Binary Cross Entropy (BCE)
      \item \textbf{VAE:} 
        \begin{itemize}
          \item Reconstruction loss (BCE)
          \item KL divergence
        \end{itemize}
      \item \textbf{Conditional VAE:}
        \begin{itemize}
          \item Reconstruction loss
          \item KL divergence
          \item Conditional generation quality
        \end{itemize}
      \item \textbf{GAN:}
        \begin{itemize}
          \item Generator loss
          \item Discriminator loss
        \end{itemize}
      \item \textbf{VQ-VAE:}
        \begin{itemize}
          \item Reconstruction loss
          \item VQ loss (commitment cost)
        \end{itemize}
      \item \textbf{Diffusion:}
        \begin{itemize}
          \item Denoising loss
          \item Sampling quality
        \end{itemize}
    \end{itemize}
  \end{block}
\end{frame}

\begin{frame}{Porównanie modeli}
  \begin{columns}[T]
    \begin{column}{0.6\textwidth}
      \begin{itemize}
        \item \textbf{Autoencoder:}
          \begin{itemize}
            \item Najszybszy trening
            \item Najprostsza implementacja
            \item Ograniczona jakość generowanych próbek
            \item BCE loss jako metryka rekonstrukcji
          \end{itemize}
        \item \textbf{VAE:}
          \begin{itemize}
            \item Lepsza jakość generowanych próbek
            \item Możliwość interpolacji w przestrzeni latentnej
            \item Stabilny trening
            \item Balans między rekonstrukcją a KL divergence
          \end{itemize}
      \end{itemize}
    \end{column}
    \begin{column}{0.4\textwidth}
      \begin{itemize}
        \item \textbf{GAN:}
          \begin{itemize}
            \item Najlepsza jakość wizualna
            \item Trudny trening
            \item Problem z mode collapse
            \item Wymaga balansu między generatorem a dyskryminatorem
          \end{itemize}
      \end{itemize}
    \end{column}
  \end{columns}
\end{frame}

\begin{frame}{Zaawansowane modele}
  \begin{columns}[T]
    \begin{column}{0.5\textwidth}
      \textbf{Conditional VAE:}
      \begin{itemize}
        \item Kontrolowane generowanie cyfr
        \item Interpolacja między cyframi
        \item Lepsza jakość niż standardowy VAE
        \item Wymaga etykiet podczas treningu
      \end{itemize}
      
      \textbf{VQ-VAE:}
      \begin{itemize}
        \item Dyskretna przestrzeń latentna
        \item Wizualizacja codebooka
        \item Dobre kompresja z zachowaniem jakości
        \item Możliwość interpolacji
      \end{itemize}
    \end{column}
    \begin{column}{0.5\textwidth}
      \textbf{Diffusion Model:}
      \begin{itemize}
        \item Progresywne generowanie
        \item Wysoka jakość próbek
        \item Długi czas generowania
        \item Stabilny proces treningu
      \end{itemize}
    \end{column}
  \end{columns}
\end{frame}

\section{Wizualizacje i przykłady}

\begin{frame}{Przykłady generowanych cyfr}
  \begin{columns}[T]
    \begin{column}{0.5\textwidth}
      \centering
      \begin{tcolorbox}[colback=yellow!10!white, colframe=black!60, boxrule=0.5pt, arc=3mm]
        \centering
        \textbf{Przykładowe wygenerowane cyfry}
        \begin{itemize}
          \item Basic AE: Rekonstrukcja cyfr
          \item VAE: Generowanie nowych cyfr
          \item GAN: Wysokiej jakości próbki
        \end{itemize}
      \end{tcolorbox}
    \end{column}
    \begin{column}{0.5\textwidth}
      \centering
      \begin{tcolorbox}[colback=yellow!10!white, colframe=black!60, boxrule=0.5pt, arc=3mm]
        \centering
        \textbf{Interpolacja w przestrzeni latentnej}
        \begin{itemize}
          \item VAE: Płynne przejścia
          \item Conditional VAE: Interpolacja między cyframi
          \item GAN: Manipulacja cechami
        \end{itemize}
      \end{tcolorbox}
    \end{column}
  \end{columns}
\end{frame}

\section{Wnioski i przyszłe kierunki}

\begin{frame}{Główne wnioski}
  \begin{block}{Najważniejsze obserwacje}
    \begin{itemize}
      \item VAE i GAN osiągnęły najlepsze wyniki w generowaniu cyfr
      \item Conditional VAE pozwala na kontrolowane generowanie
      \item Diffusion Model wymaga najdłuższego czasu treningu
      \item VQ-VAE zapewnia dobre kompresję z zachowaniem jakości
      \item Każdy model ma swoje unikalne zalety i zastosowania
    \end{itemize}
  \end{block}
\end{frame}

\begin{frame}{Kierunki rozwoju}
  \begin{columns}[T]
    \begin{column}{0.6\textwidth}
      \textbf{Możliwe ulepszenia:}
      \begin{itemize}
        \item Implementacja bardziej zaawansowanych architektur
        \item Optymalizacja hiperparametrów
        \item Eksperymenty z większymi zbiorami danych
        \item Badanie wpływu różnych funkcji straty
        \item Integracja najlepszych cech różnych modeli
      \end{itemize}
    \end{column}
    \begin{column}{0.4\textwidth}
      \begin{tcolorbox}[colback=yellow!10!white, colframe=black!60, boxrule=0.5pt, arc=3mm]
        \centering
        \textbf{Perspektywy:}
        \begin{itemize}
          \item Zastosowania praktyczne
          \item Integracja modeli
          \item Nowe architektury
        \end{itemize}
      \end{tcolorbox}
    \end{column}
  \end{columns}
\end{frame}

\begin{frame}
  \centering
  \Huge Dziękujemy za uwagę!
\end{frame}

\end{document} 